% !TeX spellcheck = it_IT
%Carattere dimensione 12
\documentclass[12pt]{report}

%Margini e interlinea
\usepackage[top=1in, bottom=1in, left=1.2in, right=1in]{geometry}
\pagestyle{plain}
\linespread{1.5}

%Librerie utili
\usepackage[italian]{babel} % conversione in italiano
\usepackage[utf8]{inputenc}
\usepackage{libertine}
\usepackage{graphicx}
\usepackage{floatflt}
\usepackage{blindtext}
\usepackage{enumitem}
\usepackage{amsthm}
\usepackage{subfig}
\usepackage{listings}
\usepackage{listingsutf8}
\usepackage{amsmath}
\usepackage{framed}
\usepackage{minibox}
\usepackage{float}
\usepackage{wrapfig}
\usepackage{longtable}
\usepackage[strict]{changepage}
\usepackage{pgfplots}
\usepackage{tikz}
\usetikzlibrary{matrix}
\pgfplotsset{width=11cm,compat=1.9}
\usepgfplotslibrary{external}
\tikzexternalize

\begin{document}

\begin{titlepage}
\begin{figure}[t]
    %\centering\includegraphics[width=0.9\textwidth]{marchio-unipi.eps}
	\centering\includegraphics[width=0.9\textwidth]{logo.eps}
	\vspace{1cm}
    \centering\includegraphics[width=0.4\textwidth]{cherubino.eps}
\end{figure}

\begin{center}
	\textbf{ Dipartimento di Informatica\\ Corso di Laurea in Informatica\\}
	\vspace{15mm}
    {\LARGE{\bf Geolocalizzazione Indoor Basata su Beacon Bluetooth a Bassa Potenza:}}\\
	\vspace{3mm}
	{\LARGE{\bf Metodi e Modelli del Deep Learning}}\\
	{{\bf un progetto realizzato per Consorzio Metis e ASL Toscana}} \\
\end{center}

\vspace{20mm}

\begin{minipage}[t]{0.47\textwidth}
	{\large{\bf Relatori:\\ Prof. Gianluigi Ferrari 
	% \\ Prof. Alessio Micheli
	}}
\end{minipage}\hfill\begin{minipage}[t]{0.47\textwidth}\raggedleft
	{\large{\bf Presentata da: \\ Marco Pampaloni \\ }}
\end{minipage}

\vspace{18mm}

\centering{\large{\bf Sessione Estiva \\ Anno Accademico 2019/2020 }}

\end{titlepage}

\tableofcontents

\begin{abstract}

\end{abstract}

\chapter{Introduzione}

\chapter{Geolocalizzazione Indoor}
\section{Introduzione al problema}
\section{Possibili soluzioni}
\section{Beacon BLE}
\section{RSSI e propagazione del segnale}
\section{Variabilità e rumore di fondo: requisiti di usabilità}

\chapter{Deep Learning}
\section{Neurone Artificiale: Perceptron}
\section{Multi Layer Perceptron}
\section{Attivazione: ReLU e Funzioni Sigmoidee}
\section{Apprendimento: Metodo del Gradiente e BackPropagation}
\section{Reti Neurali Convoluzionali}
\section{Regolarizzazione}
\subsection{Overfitting e Underfitting}
\subsection{Regolarizzazione L2}
\subsection{Dropout}
\section{Dataset Augmentation e Preprocessing}
\subsection{Jittering}
\subsection{Ridimensionamento (Scaling)}
\subsection{Magnitude Warping}
\subsection{Permutazione di Sottoinsiemi (Subset Shuffling)}
\subsection{Deattivazione Selettiva}


\chapter{Architettura Software}
\section{TensorFlow}
\section{Keras}
\section{Google Colab}
\section{Weights \& Biases}
\section{Modello di Apprendimento}
\subsection{Input del modello}
\subsection{Blocco Convoluzionale}
\subsection{Uso della Bussola e Output Ausiliario}
\subsection{Coefficiente di Memoria Residua e Input Ausiliario}
\subsection{Output del modello}
\section{Addestramento del Modello}
\section{Ensembling}
\section{Compilazione e Deploy del modello}

\chapter{Applicazione Mobile}
\section{Flutter}
\section{Planimetrie e Poligoni}
\section{Backend TensorFlow}
\subsection{TensorFlow Lite}
\subsection{Implementazione del Bridge di Comunicazione}

\chapter{Conclusioni}
\section{Risultati Sperimentali}
\subsection{Metriche di Errore: MSE, MAE, MaxAE}

\section{Lavori futuri}
\subsection{Input a Lunghezza Variabile}
\subsection{Reti Neurali Residuali}
\subsection{Variational Autoencoder: Generazione di nuovi dati}
\subsection{Transfer Learning}
\subsection{Input Masking e Ricostruzione dei Segnali}
\subsection{Transformers per Problemi di Regressione}
\subsection{Simulatore BLE}

\appendix

\bibliography{references}
\bibliographystyle{plain}

\end{document}
